\section*{Prelude}
\addcontentsline{toc}{section}{Prelude}
\label{sec:prelude}

\noindent
Some weeks before this paper was begun, its author sat down to write a betrothal document. It was written on red silk, in classical Chinese, on a day chosen because the almanac called it auspicious; and it was written in the full knowledge that nothing in the world would make it binding. No court would enforce it. No clause within it could be sued upon. It named no penalty and conferred no remedy. Against the day it might be broken it set no sword, and behind the hand that wrote it stood no power that could compel the hand to keep it. By every standard that a contract is measured against, it was a defective instrument: a promise that promised nothing it could be made to perform.

And yet it was not written lightly, and it was not, by the two for whom it was made, received lightly. The choosing of the day, the silk, the gravity of the brush---these were not theatre, or were not only theatre. Something was being done in the writing of it that the writing of a contract does not do, and that the document's want of force did not diminish but, strangely, seemed to deepen. A contract one signs and files; this was kept. The very absence of the sword behind it was felt, by the one who wrote it, not as a lack to be apologised for but as the condition of the thing's being what it was. Had there been a penalty, the promise would have been a smaller thing. That it bound nothing was the reason it could bind everything.

This is a strange experience to have, and stranger still to find that one cannot account for it. The whole apparatus by which the modern world thinks about promises---the contract, the penalty, the enforceable obligation, the institution that stands ready to compel---was built precisely to supply what this document conspicuously lacked. One is told, and half believes, that a promise is reliable in proportion to the force that backs it; that what cannot be enforced cannot be counted on; that to bind oneself is to expose oneself to a consequence one would rather avoid. By that account the betrothal document is a sentimental nullity, a pretty gesture with nothing underneath. But that is not how it was made, and it is not how it was met. Beneath the account one half believes there runs a contrary conviction, older and harder to dislodge: that the promise which asks for no sword is not the weaker but the truer; that a bond one could be \emph{forced} to keep is, just insofar as one could be forced, not yet fully a bond at all.

The author does not propose to resolve this by deciding that one of these convictions is naive. Both are real, and the tension between them is not a confusion to be cleared up but a fault line to be followed down. For the same experience that gave rise to the document gave rise, in the days after it, to a more anxious and more practical question, which is in truth the same question wearing working clothes. The promise had been made, and was believed, between the two who made it. But it had now to be carried outward---to a family, in particular, who did not yet know the author, who would not read red silk as the two who wrote it read it, and to whom the gravity of a brush on an auspicious day would say, if anything, very little. They are people, as the author understands them, who hold that what is the case is the case and what is not is not\,---\,\zh{有就是有，没有就是没有}\,---\,and who would ask, reasonably and without malice, what there is to go on. To them the bare sincerity of the vow is not evidence. And here the comfortable thought that a true promise needs no sword runs straight into its own difficulty: for if the truest promise is precisely the one that offers no external guarantee, it is also, and for the same reason, the one that gives the person who must be reassured the least to hold on to. The very purity that makes it a vow makes it, to one who asks for grounds, almost nothing at all.

So the experience opens, at its near edge and at its far, onto a single question, and it is the question this paper exists to ask. \emc{How does a promise that has no force generate trust?}\, How does an utterance about a future that does not yet exist, which threatens no consequence and can compel no compliance, produce in another person, here and now, a real expectation that it will be kept---and produce it, if the older conviction is right, more surely the less it relies on being able to compel anything at all? The question is at once the most intimate and the most public that can be asked. It is the question of how the author is to be trusted by those he has not yet earned the trust of; and it is, the paper will argue, the very question on which the law itself, in its long retreat from the sword, has been quietly working for a century. It begins in the choosing of a day and the gravity of a brush, and it reaches, before it is done, as far as the foundations of obligation and the descent of what an older language would have called the divine.

This is where the paper begins: not with a thesis but with a vow that should not have worked, and did.
